\section{宪法的基本理论}
\begin{enumerate}
  \item 制定和实施（修改、解释、监督、宣誓）、渊源和结构、规范、效力。
  \item 核心功能。分权，人权。依宪治国 = 宪政理念。
  \item 权利分配体系。
        \begin{enumerate}
          \item 立法法 \yjt 人大 \yjt 选举，监察。\zd{立法法、人大组织法、选举法}。
          \item 中央机关。中央国家机关的相关组织法。
          \item 全人大、全人常。
          \item 民族自治、特别行政区。民族区域自治法、特区基本法。
          \item 公民的基本权利、义务。
        \end{enumerate}
  \item 特征。
        \begin{enumerate}
          \item 根本法。
                \begin{enumerate}
                  \item 内容最重要。\bs{范围封闭性}。
                        \begin{enumerate}
                          \item 规定了\zd{最核心、根本}的问题。国体、政体、国家的根本制度、根本任务。
                          \item 命题，把不是宪法规定的东西说成是宪法的，\zd{选项大概率是错的}。\al{生命权 \cuo}。
                        \end{enumerate}
                  \item 效力最高。普通法的\zd{立法基础、不得违背、最高准则}。\bs{对法、对人}。
                  \item 程序最严格。修改。\bs{51 全人常，临时开会提修宪，法律主席相对多，宪法主席绝对多}。
                \end{enumerate}
          \item 权利保障书。控权保民。
          \item 民主基石。推翻专制，民主革命的产物。近代以来的产物。
        \end{enumerate}
  \item \bs{分类}。常考。分类标准。西耶士、凯尔森。
        \begin{enumerate}
          \item 有无统一的宪法典。\zd{成文、不成文}（\al{英国}）。蒲莱士。\zd{1787 美宪一、1791 法欧一}。
          \item \zd{效力、制定修改程序}是否和普通法一样。\zd{刚性、柔性}。蒲莱士。
          \item 制宪主体谁发挥主导作用。皇帝、人民、皇帝 + 人民。
                \begin{enumerate}
                  \item 钦定。\al{1889 日本明治宪法、1908 钦定宪法大纲（中国第一个宪法文件）}。
                  \item 民定。\al{xx 共和国}。
                  \item 协定。\al{1215 英国自由大宪章（世界第一个宪法文件）、1803 法国宪法}。
                \end{enumerate}
        \end{enumerate}
  \item \zd{4 个基本原则}。\zd{人民}主权（首要）、基本\zd{人权}、\zd{法治}原则（良法之治 + 普遍守法）、\zd{权力}制约。
  \item 制定。\zd{人民制宪（本源）、人大修宪、人常解释宪法、人大加人常监督宪法的实施、78 给解释、82 给实施。
          人大制定基本法，人常制定基本法律以外的普通法律。
          人大可以对人常进行基本法的授权。
          人常可以在人大闭会期间，在不违反基本法律原则和精神的前提下对基本法律作出部分的修改和补充。
          但是，两个香港、澳门基本法除外。
          并且人常会还可以进行立法解释和立法裁决来解释法律内部的冲突}。
        \begin{enumerate}
          \item 设立制宪机关。按专门性又分为\zd{起草}、通过机关。
          \item 提出宪法草案。全民讨论。
          \item 通过。通过机关 \hz 全民公决。不可能有具体数字。
          \item 公布。元首模式 \hz \zd{议会模式}。我国\zd{全人大发布公告 = 主席团发布人大公告}。
          \item 制宪权、制宪机关一次性的。
        \end{enumerate}
  \item 我国宪法历史。5 次修改，52 条修正案。3 大 7 小 1 制定。必考。
        \begin{enumerate}
          \item 近代宪法、现代宪法。法制史内容。
          \item 共同纲领（过渡，新民主主义） \yjt 54 宪法 \yjt 75 \yjt 78 \yjt 82。
          \item \zd{修改方式。人大决议、修正案}（82 开始用，第一次使用是 88）。
          \item 52 条修正案。
                \begin{enumerate}
                  \item 88。邓。私营。
                  \item 93。邓。国营改国有、市场经济、初级阶段、多党合作政治协商。
                  \item 99。江。依法治国、反革命改危害国家安全、个体私营改非公有、公有制为主体。
                  \item 04。胡。\zd{人}权、\zd{紧}急状态、\zd{社}会保障、\zd{建}设者入\zd{站}；增加\zd{国歌}、\zd{国事}访问、\zd{特区}人大、\zd{政治}文明。
                  \item 18。习。社会生态、和谐美丽、伟大复兴、科学发展、习近平。
                  \item \bs{民营经济。88 补充，93 市场，99 重要，04 加上鼓励和支持}。
                  \item \bs{土地。88 转，93 分，99 统，04 征}。
                  \item \bs{任期。93 县，04 乡，3 改 5}。
                  \item \bs{04 建设，18 复兴}。
                \end{enumerate}
        \end{enumerate}
  \item 宪法渊源。必考。
        \begin{enumerate}
          \item 宪法典。我国 1 托 5，1 部宪法典，5 个修正案。
          \item 宪法性法律。宪法的部门法。组织法、选举法、代表法、立法法、代议机关的议事规则。
          \item 宪法惯例。没有国家强制力而已。
          \item 宪法判例。美国、英国。我国没有，只能援引，不能拿宪法作为判决依据。\zd{修正案可推翻判例}。
          \item 国际条约。和平共处五项原则 \yjt 和平、共赢、命运。
        \end{enumerate}
  \item 宪法典结构。序言（推定有法律效力）、正文、附则（\yi 确定\zd{有效}。\er \bs{我国没有附则}）。
  \item 宪法规范。
        \begin{enumerate}
          \item 特征。根本性、最高性、原则性、纲领性、稳定性。
          \item 分类。
                \begin{enumerate}
                  \item 确认性、禁止性、强制性规范。
                  \item 权力性、义务性、权利义务复合。
                  \item 程序。直接、间接。
                \end{enumerate}
        \end{enumerate}
  \item 效力。正当性、最高性、直接性。及于所有的中国公民和华侨。
  \item 趋势。国家权利扩大，行政权力受到一定程度限制。
\end{enumerate}
